% sec5_4.tex — Section 5.4: Application: International Differences in Growth Rates

\section{Application: International Differences in Growth Rates}
\label{sec:5.4}

Do poor economies grow faster than rich economies? If so, how much faster?
These are the questions addressed in the recent burgeoning literature of growth.\footnote{A
very extensive survey of the literature is Barro and Sala-i-Martin (1995).}
This section derives the key equation that has been used to explain economic
growth of nations and discusses ways to estimate it. Actual estimation of several
specifications of the equation is relegated to the empirical exercise.

\subsection*{Derivation of the Estimation Equation}

In neoclassical growth theory, output per effective labor in period $t$, denoted $q(t)$
and to be defined more precisely in \eqref{eq:5.4.4} below, converges to the steady-state
level $q^*$. The log-linear approximation around the steady state gives
\begin{equation}
  \frac{\mathrm{d}\log(q(t))}{\mathrm{d}t}
  = \lambda \cdot [\log(q^*) - \log(q(t))].
  \label{eq:5.4.1} \tag{5.4.1}
\end{equation}

Estimating $\lambda$, the \textbf{speed of convergence}, has been the objective of the recent
growth literature. \eqref{eq:5.4.1} implies that, for any two points, $t_{m-1}$ and $t_m$ in time,
\begin{equation}
  \log(q(t_m)) = (1 - \rho) \cdot \log(q^*) + \rho \cdot \log(q(t_{m-1})),
  \label{eq:5.4.2} \tag{5.4.2}
\end{equation}
where
\begin{equation}
  \rho \equiv \exp[-\lambda \cdot (t_m - t_{m-1})].
  \label{eq:5.4.3} \tag{5.4.3}
\end{equation}

(Since in our data $t_m$'s are equally spaced in time, $\rho$ will not depend on $m$.)
Output per effective labor is defined as
\begin{equation}
  q(t) \equiv \frac{Y(t)}{A(t)\,L(t)},
  \label{eq:5.4.4} \tag{5.4.4}
\end{equation}
where $Y(t)$ equals aggregate output, $L(t)$ equals aggregate hours worked, and $A(t)$
equals level of labor-augmenting technical progress. Assuming $A(t)$ grows at a
constant rate $g$ (so that $A(t) = A(0)\exp(gt)$), \eqref{eq:5.4.4} implies
\begin{equation}
  \log(q(t)) = \log\biggl(\frac{Y(t)}{L(t)}\biggr) - \log(A(0)) - g \cdot t.
  \label{eq:5.4.5} \tag{5.4.5}
\end{equation}

Substituting this into \eqref{eq:5.4.2}, we obtain
\begin{equation}
  \log\biggl(\frac{Y(t_m)}{L(t_m)}\biggr)
  = \rho \cdot \log\biggl(\frac{Y(t_{m-1})}{L(t_{m-1})}\biggr)
  + (1-\rho) \cdot [\log(q^*) + \log(A(0))] + \phi_m,
  \label{eq:5.4.6} \tag{5.4.6}
\end{equation}
where
\begin{equation}
  \phi_m \equiv g \cdot (t_m - \rho \cdot t_{m-1}).
  \label{eq:5.4.7} \tag{5.4.7}
\end{equation}

An equivalent form of equation \eqref{eq:5.4.6} can be obtained by subtracting
$\log\frac{Y(t_{m-1})}{L(t_{m-1})}$ from both sides:
\begin{align}
  \log\biggl(\frac{Y(t_m)}{L(t_m)}\biggr) - \log\biggl(\frac{Y(t_{m-1})}{L(t_{m-1})}\biggr)
  &= (\rho - 1)\log\biggl(\frac{Y(t_{m-1})}{L(t_{m-1})}\biggr)
  \notag \\
  &\quad + (1-\rho)\cdot[\log(q^*) + \log(A(0))] + \phi_m.
  \label{eq:5.4.6p} \tag{5.4.6$'$}
\end{align}

Since $\rho < 1$ as $\lambda > 0$, this equation says that the level of per capita output has a
negative effect on subsequent growth. The country, when poor, should grow faster.
This form of the equation is more intuitive, but for the purpose of choosing the
correct estimation technique, \eqref{eq:5.4.6} rather than \eqref{eq:5.4.6p} is more instructive.

Equation \eqref{eq:5.4.6} should hold for each country $i$.\footnote{We have derived the equation
describing convergence assuming a closed economy. The assumption is hard
to justify, but a similar equation can be derived under free mobility of physical capital
if human capital is another
factor of production and not freely mobile internationally. See, e.g., Barro and
Sala-i-Martin (1995).}
For the rest of this section, we
assume that the speed of convergence ($\lambda$) and the growth rate of labor-augmenting
technical progress ($g$) are the same across countries. Then \eqref{eq:5.4.6} for country $i$ can
be written as
\begin{equation}
  y_{im} = \phi_m + \rho\, y_{i,m-1} + \alpha_i,
  \label{eq:5.4.8} \tag{5.4.8}
\end{equation}
where
\begin{align}
  y_{im} &= \log\biggl(\frac{Y(t_m)}{L(t_m)}\biggr) \quad \text{for country } i,
  \notag \\
  \alpha_i &= (1-\rho) \times \{\log(q^*) + \log(A(0))\}
  \quad \text{for country } i.
  \label{eq:5.4.9} \tag{5.4.9}
\end{align}

\subsection*{Appending the Error Term}

It is a fairly standard practice in applied econometrics to derive an equation without
an error term from the theory and then append an error term for estimation. Growth
empirics is not an exception; an error term $\eta_{im}$ is added to \eqref{eq:5.4.8} to obtain
\begin{equation}
  y_{im} = \phi_m + \rho\, y_{i,m-1} + \alpha_i + \eta_{im}.
  \label{eq:5.4.10} \tag{5.4.10}
\end{equation}

One natural interpretation of $\eta_{im}$ is business cycles, which in macroeconomics are
defined as the difference between actual output and potential output. Growth theory
is about potential output. The discrepancy between potential output and actual
output becomes the error term in the estimation equation derived from growth theory.
But then there would be an errors-in-variables problem, which renders
the regressors endogenous.\footnote{Let $x_{im}$ be the log potential output in year $t_m$
for country $i$ that growth theory is purported to explain, so that
(5.4.8) holds for $x_{im}$ rather than for $y_{im}$. Let log actual output be related to
$x_{im}$ as $y_{im} = x_{im} + \eta_{im}$. Then
(5.4.10) results, but with the error term $\eta_{im} - \rho\eta_{i,m-1}$.
Not only would there be an errors-in-variables problem,
but also the error term has a moving-average structure.}
For the rest of this section, we follow the mainstream
growth literature and ignore the endogeneity problem.

Another problem is the possible correlation in $\eta_{im}$ between countries ($i$'s). This
spatial (or geographical) correlation would not arise if the sample were a random
sample. But the sample of nations is not a random sample; it is the universe.
The phenomenon known as international business cycles implies that the spatial
correlation would be positive. Statistical inference treating countries as if they were
independent data points would overstate statistical significance. This problem, too,
is largely ignored in the literature and will be ignored in our discussion.

\subsection*{Treatment of $\alpha_i$}

As is clear from \eqref{eq:5.4.9}, the term $\alpha_i$ in equation \eqref{eq:5.4.10} depends on the steady-state
level $q^*$ and the initial level of technology $A(0)$. How you treat $\alpha_i$ in the
estimation depends on which version of neoclassical growth theory you subscribe
to. According to the Cass-Koopmans optimal growth model, the sole determinant
of the steady-state level $q^*$ is the growth rate of labor. Thus, if technological
knowledge freely flows internationally so that $A(0)$ is the same across countries,
the international differences in $\alpha_i$ are completely explained by the labor growth
rate. Therefore, leaving aside the endogeneity problem just mentioned, OLS estimation
of \eqref{eq:5.4.10} delivers consistency if the labor growth rate is included as the
additional regressor to control for $\alpha_i$. If you subscribe to the Solow-Swan growth
theory, the saving rate is the other determinant of $q^*$, so the regression should
include the saving rate along with the labor growth rate.

A large fraction of the recent growth literature under the rubric of ``conditional
convergence'' includes these and other variables that might affect $q^*$ or $A(0)$ ---
such as measures of political stability and the degree of financial intermediation ---
in order to control for $\alpha_i$. In the empirical exercise, you will be asked to estimate
$\rho$, and hence the speed of convergence $\lambda$, using this conditional convergence
approach.

However, no matter how ingeniously the variables are selected to control for
$\alpha_i$, some component of $\alpha_i$ having to do with features unique to the country would
remain unexplained and find its way in the error term. Being part of $\alpha_i$, it is correlated
with the regressor $y_{i,m-1}$ in \eqref{eq:5.4.10}. The OLS estimation, even if the equation
includes a variety of variables to control for $\alpha_i$, will not provide consistency. The
alternative approach is to treat $\alpha_i$ as an unobservable fixed effect.

\subsection*{Consistent Estimation of Speed of Convergence}

If output and other variables are observed at $M + 1$ points in time,
$t_0, t_1, t_2, \ldots, t_M$,
equation \eqref{eq:5.4.10} is available for $m = 1, \ldots, M$. The fixed-effects technique could
be applied to this $M$-equation system, but it does not provide consistency. The reason
is that the system is ``dynamic'' in that some of the regressors are the dependent
variables for other equations of the system. To see this, look at the first two equations
of the system:
\begin{align*}
  y_{i1} &= \phi_1 + \rho\, y_{i0} + \alpha_i + \eta_{i1}, \\
  y_{i2} &= \phi_2 + \rho\, y_{i1} + \alpha_i + \eta_{i2}.
\end{align*}

Suppose the error term $\eta_{i1}$ is orthogonal to the regressors (which are a constant and
$y_{i0}$) in the first equation, so that $\E(\eta_{i1}) = 0$ and $\E(y_{i0}\,\eta_{i1}) = 0$.
If we further make
the assumption that $\E(\alpha_i\,\eta_{i1}) = 0$, then from the first equation
\[
  \E(y_{i1}\,\eta_{i1}) = \Var(\eta_{i1}) \neq 0.
\]

But $y_{i1}$ is one of the regressors in the second equation. Therefore, if the error term
$\eta_{im}$ is orthogonal to the regressors for each equation, the ``cross'' orthogonalities are
necessarily violated.\footnote{This was first pointed out by Nickell (1981).}
As emphasized in Section~5.2 (see Review Question~3), the
fixed-effects estimator is not guaranteed to be consistent if the cross orthogonalities
are not satisfied. For a precise expression for the bias, see Analytical Question~4 to
this chapter.

One way to obtain consistency is the following. Take the first differences of
the $M$-equations to obtain $M - 1$ equations:
\begin{align}
  y_{i2} - y_{i1} &= \mu_1 + (y_{i1} - y_{i0})\rho + (\eta_{i2} - \eta_{i1}),
  \notag \\
  y_{i3} - y_{i2} &= \mu_2 + (y_{i2} - y_{i1})\rho + (\eta_{i3} - \eta_{i2}),
  \notag \\
  &\;\;\vdots \notag \\
  y_{iM} - y_{i,M-1} &= \mu_{M-1} + (y_{i,M-1} - y_{i,M-2})\rho
  + (\eta_{iM} - \eta_{i,M-1}),
  \label{eq:5.4.11} \tag{5.4.11}
\end{align}
where
\[
  \mu_m \equiv \phi_{m+1} - \phi_m.
\]

The random-effects estimator and the equation-by-equation OLS do not provide
consistency because the regressors are not orthogonal to the error term. For example,
in the first equation, $y_{i1} - y_{i0}$ is not orthogonal to $\eta_{i2} - \eta_{i1}$
because $\E(y_{i1} \cdot \eta_{i1}) \neq 0$ as just seen. Consistent estimation can be done by
multiple-equation GMM, if
there are valid instruments. In the optional empirical exercise to this chapter, you
will be asked to use the saving rate as the instrument for the endogenous regressors.
To the extent that the saving rate is uncorrelated with $\eta_{im}$, this estimator is robust
to the endogeneity problem mentioned above.

\bigskip
\noindent\rule{\textwidth}{0.5pt}
\textsc{Questions for Review}
\medskip

\begin{enumerate}
\item The assumption that $g$ (growth rate of technical progress) is common to all
  countries is needed to derive \eqref{eq:5.4.8}. Do we have to assume that $g$ is constant
  over time? [Answer: No.]

\item (Identification)\quad For simplicity, let $M = 3$ so that \eqref{eq:5.4.11} is a two-equation
  system.
  \begin{enumerate}[label=(\alph*)]
  \item Write \eqref{eq:5.4.11} as a multiple-equation system with common coefficients by
    properly defining the $y_{i1}$, $y_{i2}$, $\mathbf{z}_{i1}$, $\mathbf{z}_{i2}$,
    $\boldsymbol{\delta}$ in (4.6.1). \textbf{Hint:}
    $\mathbf{z}_{i2}$, for example, is $(0, 1, y_{i2} - y_{i1})'$.
  \item If $s_{im}$ is the instrument for $y_{im} - y_{i,m-1}$, the instrument vector for
    the $m$-th equation is $\mathbf{x}_{im} = (1, s_{im})'$ for $m = 1, 2$. Show that the system is
    not identified if $\Cov(s_{im},\, y_{im} - y_{i,m-1}) = 0$ for all $m$.
    \textbf{Hint:}
    \[
      \boldsymbol{\Sigma}_{\mathbf{xz}} = \begin{bmatrix}
        1 & 0 & \E(y_{i1} - y_{i0}) \\
        \E(s_{i1}) & 0 & \E[s_{i1}(y_{i1} - y_{i0})] \\
        0 & 1 & \E(y_{i2} - y_{i1}) \\
        0 & \E(s_{i1}) & \E[s_{i2}(y_{i2} - y_{i1})]
      \end{bmatrix}.
    \]
  \end{enumerate}
\end{enumerate}
